\section*{第 \ref{ch:g12:innate-immunity} 章——先天免疫}

\begin{solution}{exo:g12:innate-immunity:1}
\omterm{def:g12:innate-immunity:immunity}{先天免疫}：出生时即已存在，数小时内起作用，识别同类入侵者，
依靠共有的分子模式，每次应答都相同。\omterm{def:g12:innate-immunity:immunity}{适应性免疫}：需要数天，
逐一识别每个入侵者并对其形成记忆（且仅存在于
\omterm{def:g10:common-ancestry:bodyplan}{脊椎动物}中）。
\end{solution}

\begin{solution}{exo:g12:innate-immunity:2}
发红与发热：血管扩张，输送更多血液。肿胀：血管壁渗漏，
使血浆进入组织。疼痛：神经末梢受\omterm{prop:g12:innate-immunity:inflammation}{介质}
致敏。
\end{solution}

\begin{solution}{exo:g12:innate-immunity:3}
同类微生物共有的分子模式（如\omterm{def:g10:cells-common-unit:organelle}{细胞壁}片段、
病毒\omterm{def:g10:chemistry-of-life:families}{核酸}），以及受损\omterm{def:g10:cells-common-unit:cell}{细胞}释放的分子。它们
促使\omterm{def:g10:cells-common-unit:cell}{细胞}释放\omterm{prop:g12:innate-immunity:inflammation}{介质}：组胺、前列腺素、\omterm{def:g10:cells-common-unit:cell}{细胞}因子。
\end{solution}

\begin{solution}{exo:g12:innate-immunity:4}
通过受体接触并结合；由\omterm{def:g10:cells-common-unit:cell}{细胞}膜包裹吞入；
在囊泡中借助酶和活性化学物质进行消化；
将\omterm{def:g11:gene-expression:protein}{蛋白质}片段呈递于\omterm{def:g10:cells-common-unit:cell}{细胞}表面。
\end{solution}

\begin{solution}{exo:g12:innate-immunity:5}
它们阻断合成前列腺素的酶，从而减轻肿胀、
疼痛和发热。它们既不清除病因，也不杀灭任何微生物。
\end{solution}

\begin{solution}{exo:g12:innate-immunity:6}
\omterm{def:g12:innate-immunity:immunity}{先天免疫}约在半天时起作用；\omterm{def:g12:innate-immunity:immunity}{适应性免疫}约在第 9 天起作用；两条曲线
约在第 5--6 天相交。
\end{solution}

\begin{solution}{exo:g12:innate-immunity:7}
触发因素是哨兵\omterm{def:g10:cells-common-unit:cell}{细胞}识别的分子模式，而非活的入侵者。
无菌生理盐水不含任何模式，也不造成损伤，
至多只会引起针刺后的轻微反应。
\end{solution}

\begin{solution}{exo:g12:innate-immunity:8}
死亡的中性粒\omterm{def:g10:cells-common-unit:cell}{细胞}、死亡和存活的细菌、被消化的碎屑以及
血浆。中性粒\omterm{def:g10:cells-common-unit:cell}{细胞}在数小时内陆续到达，并在约一天内死亡；
脓液便是它们战斗后积累的产物。
\end{solution}

\begin{solution}{exo:g12:innate-immunity:9}
花粉热是肥大\omterm{def:g10:cells-common-unit:cell}{细胞}针对花粉释放组胺所致：阻断组胺即可消除该反应。
在感染中，组胺只在最初几分钟起作用；
之后的肿胀、疼痛和\omterm{def:g10:cells-common-unit:cell}{细胞}募集依赖前列腺素和\omterm{def:g10:cells-common-unit:cell}{细胞}因子，
抗组胺药对它们并无作用。
\end{solution}

\begin{solution}{exo:g12:innate-immunity:10}
在 12 小时时，分别为 100\% 对 55\%；在 36 小时时，分别为
50\% 对 30\%。该药物在各时间点均能降低肿胀，但曲线仍在约
72 小时降至零：它并未缩短应答的时间。
\end{solution}

\begin{solution}{exo:g12:innate-immunity:11}
撕裂的\omterm{def:g10:muscles-and-joints:joint}{韧带}纤维和被压碎的\omterm{def:g10:cells-common-unit:cell}{细胞}释放出受损\omterm{def:g10:cells-common-unit:cell}{细胞}的分子，
哨兵\omterm{def:g10:cells-common-unit:cell}{细胞}的受体像识别微生物一样识别这些分子：
仅损伤本身就足以触发应答。
\end{solution}

\begin{solution}{exo:g12:innate-immunity:12}
没有中性粒\omterm{def:g10:cells-common-unit:cell}{细胞}，最初几天便会失守，普通细菌将压垮机体：
先天应答的作用正在于立即守住防线。没有\omterm{def:g12:innate-immunity:immunity}{适应性免疫}时，
先天应答能清除最初的感染，但持续存在或反复发生的感染，
需要特异性抗体和记忆才能应对，最终可能占据上风：
适应性应答所完成的，正是先天应答所遏制下来的
那部分任务。
\end{solution}

\begin{solution}{exo:g12:innate-immunity:13}
脓肿是针对活细菌的炎症：抑制它会减慢细菌的清除，
并有扩散的风险。\omterm{def:g10:muscles-and-joints:joint}{关节}炎是没有任何微生物参与的炎症，
直接针对\omterm{def:g10:muscles-and-joints:joint}{关节}本身：此时应答即是疾病，
减轻应答便是治疗。
\end{solution}

\begin{solution}{exo:g12:innate-immunity:14}
先天受体由\omterm{def:g10:universal-dna:gene}{基因}决定，识别整类生物共有的分子模式；
再次相遇并不会改变它们。适应性系统则为每个入侵者
构建特异性受体，并保留产生这些受体的\omterm{def:g10:cells-common-unit:cell}{细胞}：
它识别的是个体，
并保留这种识别能力。
\end{solution}

\begin{solution}{exo:g12:innate-immunity:15}
炎症针对真正的入侵者或损伤、程度适当，并在数天内消退时，
是有益的：它清除病灶并启动修复。
当它针对无害之物（如花粉、机体自身的\omterm{def:g10:muscles-and-joints:joint}{关节}）、过度，
或持续不断时，其\omterm{prop:g12:innate-immunity:inflammation}{介质}和\omterm{def:g12:innate-immunity:phagocytosis}{吞噬细胞}便会损伤它本应保护的组织，
炎症便成为疾病。决定因素在于靶标和持续时间，
而非机制。两种情形下，
机制是相同的。
\end{solution}

\begin{solution}{pb:g12:innate-immunity:1}
\textbf{1.} 2 小时内倍增三次：$\num{10000} \times 8 =
\num{80000}$；4 小时内倍增六次：$\num{640000}$。

\textbf{2.} 每小时 $200 \times 5 = 1000$ 个——远低于第一小时
倍增所新增的 $\num{10000}$ 个细菌，
因此增长仍在继续。

\textbf{3.} $\num{50000} \times 20 = 10^6$ 个细菌。

\textbf{4.} 第 2 小时时细菌数少于 $\num{80000}$；仅第一小时的
中性粒\omterm{def:g10:cells-common-unit:cell}{细胞}就能清除一百万个。从第 2 小时起，\omterm{def:g12:innate-immunity:phagocytosis}{吞噬细胞}
开始占据优势。

\textbf{5.} 数小时内细菌便已消失；剩下的是死亡的
中性粒\omterm{def:g10:cells-common-unit:cell}{细胞}、碎片和液体，即将形成脓液，而巨噬\omterm{def:g10:cells-common-unit:cell}{细胞}正在
将其清除。

\textbf{6.} 发红和发热：组胺和前列腺素使血管扩张。
肿胀：组胺使血管壁渗漏，血浆渗出至组织中。疼痛：
前列腺素使神经末梢致敏。

\textbf{7.} 血管渗漏使中性粒\omterm{def:g10:cells-common-unit:cell}{细胞}和血浆中的防御蛋白
能够到达组织；压迫已致敏神经的液体，
便是疼痛的原因。

\textbf{8.} 每次心跳都会将更多血液泵入扩张的血管；
肿胀的组织无法扩张，因此每次搏动都会提高对已致敏神经
的压力，形成与心跳同步的抽痛。

\textbf{9.} 死亡的中性粒\omterm{def:g10:cells-common-unit:cell}{细胞}、死亡和存活的细菌、
被消化的碎片、血浆。

\textbf{10.} 随着细菌消失，碎片被清除，模式
也随之消失：哨兵\omterm{def:g10:cells-common-unit:cell}{细胞}不再识别，也没有\omterm{prop:g12:innate-immunity:inflammation}{介质}
产生；血管恢复正常。

\textbf{11.} 它能减轻由组胺引起的早期发红、发热和渗漏；
但不能减轻由前列腺素引起的后期肿胀和疼痛，也不能
阻止中性粒\omterm{def:g10:cells-common-unit:cell}{细胞}的募集。

\textbf{12.} 它抑制合成前列腺素的酶：疼痛、部分肿胀和
发热随之消退。

\textbf{13.} 该药膏可关闭大部分\omterm{prop:g12:innate-immunity:inflammation}{介质}，因此招募的
中性粒\omterm{def:g10:cells-common-unit:cell}{细胞}减少；\omterm{def:g12:innate-immunity:phagocytosis}{吞噬细胞}减少后，细菌可能无法被清除，
感染可能扩散。

\textbf{14.} 这三种药物都作用于症状，但都不能杀灭细菌；
皮质类固醇通过减少\omterm{def:g12:innate-immunity:phagocytosis}{吞噬细胞}的募集，
可能使感染持续更久。

\textbf{15.} \omterm{def:g10:cells-common-unit:cell}{细胞}因子到达大脑，提高体温调定点；
额外的热量加快\omterm{def:g12:innate-immunity:phagocytosis}{吞噬细胞}的作用，同时减缓细菌的生长。

\textbf{16.} 它们在感染部位摄取的\omterm{def:g10:cells-common-unit:cell}{细胞}因子，告知
淋巴\omterm{def:g10:cells-common-unit:cell}{细胞}入侵者的类型；随后它们将信息传递给
淋巴结中的淋巴\omterm{def:g10:cells-common-unit:cell}{细胞}。

\textbf{17.} 大约一周至十天。对于被\omterm{def:g12:innate-immunity:immunity}{先天免疫}在数小时内
清除的木刺，它们并非必需——但仍会被产生，
并被记住。

\textbf{18.} 得益于记忆，\omterm{def:g12:innate-immunity:immunity}{适应性免疫}在第二次应答时更快、
更强；\omterm{def:g12:innate-immunity:immunity}{先天免疫}则完全相同。

\textbf{19.} \omterm{def:g11:antibiotic-resistance:antibiotic}{抗生素}可杀灭细菌，在细菌感染中替代缺失的\omterm{def:g12:innate-immunity:phagocytosis}{吞噬细胞}；
但它们无法替代\omterm{def:g12:innate-immunity:phagocytosis}{吞噬细胞}清除碎片、
对真菌的作用，
或向淋巴结传递信息。

\textbf{20.} 约在第 2 小时，\omterm{def:g12:innate-immunity:phagocytosis}{吞噬细胞}赶上细菌；
组胺和前列腺素分别导致发红、发热和肿胀，前列腺素导致疼痛；
\omterm{prop:g12:innate-immunity:handover}{树突状细胞}负责传递信息。
\end{solution}
